    \chapter{Conclusões e Trabalho Futuro}

    \section{Conclusão}

    No momento em que o desenvolvimento deste projeto teve início, a adoção de componentes de infraestrutura dedicados a intermediar aplicações e APIs de inferência era ainda uma prática pouco comum. A falta de arquiteturas e de produtos consolidados fazia com que a maioria das aplicações integrasse diretamente as APIs dos fornecedores, o que limitava a escalabilidade, a segurança e a gestão de chaves de API, de utilizadores e de outras questões fundamentais, como a imposição de limites, a implementação de mecanismos de resiliência e a recolha de métricas.

    Contudo, à medida que este projeto foi sendo desenvolvido, começaram a surgir cada vez mais soluções de intermediação de inferência (\textit{AI Gateways}), reforçando ainda mais a importância do problema que procurámos resolver. Esta evolução veio demonstrar que o projeto desenvolvido é relevante, não apenas pelos desafios técnicos que apresenta, mas também pela crescente valorização deste tipo de soluções no mercado.

    A principal diferença introduzida por este projeto é o facto de ter sido desenvolvido um \textit{Software Development Kit} (\textbf{SDK}) e não apenas uma aplicação final. Foi desenvolvida uma biblioteca completamente abstrata e customizável, recorrendo a \textbf{Kotlin Multiplatform}. Esta abordagem torna-se particularmente interessante porque não disponibiliza apenas um produto pronto a utilizar, com regras de negócio específicas, mas sim uma base reutilizável sobre a qual qualquer utilizador pode construir a solução que melhor se adapta às suas necessidades. Apesar de esta decisão ter aumentado a complexidade do projeto, acreditamos que foi também um dos aspetos que mais o diferencia das soluções existentes.

    O uso de \textbf{Kotlin Multiplatform} permite instanciar uma \textit{OmniAiGateway} em dois \textit{targets} distintos, JVM e Node.js. Embora tenham existido dificuldades associadas ao desenvolvimento de uma biblioteca em \textbf{Kotlin Multiplatform}, estas foram ultrapassadas, permitindo disponibilizar uma solução em ambas as plataformas. Este aspeto aumenta ainda mais a flexibilidade do projeto, facilitando a sua adoção em sistemas já existentes, quer na JVM quer em Node.js.

    Ao longo deste trabalho foi possível concretizar um \textbf{SDK} reutilizável e abstrato, incorporando as funcionalidades apresentadas ao longo deste relatório. Foi possível normalizar as APIs de inferência da \textbf{OpenAI} \cite{openai_api}, \textbf{Anthropic} \cite{anthropic_api} e \textbf{Gemini}, suportando igualmente rotas de entrada compatíveis com as mesmas e permitindo encaminhar os pedidos para APIs de inferência compatíveis com esses formatos. Isto permite que aplicações de IA generativa já existentes, desenvolvidas para qualquer um destes fornecedores, consigam utilizar o nosso sistema praticamente sem alterações, mantendo compatibilidade com os formatos esperados enquanto passam a beneficiar de funcionalidades adicionais.

    Sendo este um projeto bastante abstrato e configuravel, é possível configurar uma \textit{gateway} com diferentes níveis de complexidade. Pode ser utilizada apenas como uma peça de infraestrutura mais simples que encaminha pedidos para o mesmo fornecedor de inferência do pedido, acrescentando funcionalidades como métricas, \textit{rate limiting}, autenticação ou autorização. Por outro lado, pode também suportar cenários muito mais avançados, como roteamento inteligente entre modelos e fornecedores com base em latência ou algoritmos semânticos, mecanismos de \textit{caching} de respostas de inferência ou mesmo ao nível do MCP. Ao optar por esta abordagem, o nosso objetivo foi precisamente permitir que quem utiliza o \textbf{SDK} tenha liberdade para implementar funcionalidades que vão além daquelas que nós próprios desenvolvemos.

    Naturalmente, este projeto apresentou também vários desafios. Um dos principais foi lidar com as diferenças existentes entre os vários fornecedores de inferência. Apesar de, à primeira vista, as APIs parecerem semelhantes, são extremamente extensas, contêm inúmeros conceitos proprietários e nem sempre possuem documentação suficientemente explícita. Construir um domínio comum sobre estas APIs foi, para nós, um dos maiores feitos alcançados neste projeto e um resultado do qual nos orgulhamos bastante.

    Outro desafio importante foi garantir a evolução e manutenção da normalização criada. As APIs de inferência encontram-se em constante evolução e, durante o desenvolvimento deste projeto, surgiram diversas atualizações que obrigaram à adaptação da nossa implementação. Além disso, normalizar não apenas pedidos e respostas, mas também eventos de \textit{streaming}, exigiu um esforço adicional devido às diferenças existentes entre fornecedores, tornando esta componente particularmente desafiante.

    De forma a demonstrar a utilização prática do \textbf{SDK}, foi realizada a instanciação de uma \textit{OmniAi Gateway} e a sua integração com o \textbf{Claude Code}. Conseguir utilizar um agente amplamente utilizado e, através da nossa \textit{gateway}, comunicar com modelos da Google, nomeadamente o \textbf{Gemini}, foi um resultado que consideramos particularmente marcante. Durante o desenvolvimento deste trabalho não encontrámos nenhuma solução que permitisse este tipo de interoperabilidade da mesma forma, pelo que acreditamos que este projeto apresenta um caráter bastante inovador.

    A integração da camada de inferência com o \textbf{MCP Broker} foi também um elemento fundamental do projeto. Para além de nos permitir aprofundar o conhecimento sobre o protocolo \textbf{MCP}, possibilitou igualmente a implementação de um sistema interessante e totalmente integrado com a restante arquitetura.

    Consideramos que, no final, a implementação do projeto apresenta uma qualidade técnica muito elevada e cumpre os objetivos inicialmente definidos. Ao longo do desenvolvimento tivemos oportunidade de aprofundar conhecimentos em arquiteturas como \textit{Ports and Adapters}, na construção de um sistema que atua como \textit{Resource Server}, implementando os RFCs apropriados, na utilização de padrões de autorização como \textit{Policy Enforcement Point} e \textit{Policy Decision Point}, bem como na adoção do \textbf{OpenTelemetry} e na construção de \textit{dashboards} no \textbf{Grafana}. O sistema continua em evolução e existem ainda muitas funcionalidades que gostaríamos de implementar. Tratando-se de um \textbf{SDK}, existe sempre espaço para acrescentar novas capacidades. Para além dos desafios técnicos, uma das maiores dificuldades foi também definir um âmbito que fosse simultaneamente ambicioso e realista. No final, acreditamos que este projeto resultou numa solução sólida e tecnicamente diferenciadora, contribuindo também de forma significativa para a evolução das nossas competências enquanto engenheiros.



    \section{Melhorias Futuras e Escalabilidade}
    Como trabalho futuro, existem várias linhas de evolução. A primeira é consolidar a distribuição pública do \textbf{SDK} via \textbf{Maven} e \textbf{NPM}, garantindo documentação de instalação, versionamento semântico e exemplos de consumo. A segunda é completar a paridade entre \textit{targets} \textbf{JVM} e \textbf{JavaScript}, sobretudo na montagem de \textit{runtime} e transporte \textbf{HTTP}.

    Outra melhoria importante é aprofundar o \textbf{MCP Broker}, incluindo carregamento efetivo de ferramentas via \textbf{YAML}, execução completa dos \textit{handlers} \textbf{MCP} e suporte robusto a \textbf{SSE} e \textbf{WebSocket}. Também será relevante melhorar as políticas de roteamento, permitindo decisões com base em custo, latência, disponibilidade, qualidade do modelo, plano do cliente e limites de utilização.

    Por fim, a camada de segurança pode evoluir para políticas de autorização mais expressivas, integração com PDPs externos, introspeção real de \textit{tokens} opacos e mapeamento mais fino entre \textit{claims}, planos e quotas. A nível de infraestrutura, perspetiva-se a adoção de um cofre de segredos (\textit{Key Vault}) \cite{ref_keyVault}  para o armazenamento e gestão centralizada e segura das \textit{API Keys} de inferência. 
    
    Na parte da testabilidade também será fundamental alargar a validação do sistema através da implementação de testes de integração mais abrangentes e de (\textit{stress testing}), garantindo a estabilidade em casos de cenários de uso intensivo. 
    
    Estas melhorias permitiriam aproximar a \textbf{OmniAI Gateway} de uma solução robusta e pronta para ambientes \textit{multi-tenant} \cite{ref_MultiTenant} e de produção.

